\chapter{Wernicke encephalopathy}
\index{Wernicke encephalopathy}

\section*{Definition and Overview}
Wernicke encephalopathy (WE) is an acute, life-threatening neuropsychiatric syndrome caused by a severe deficiency of thiamine (vitamin B1). First described by the German neurologist Carl Wernicke in 1881, the condition historically presented as a clinical triad of symptoms: acute mental confusion, ophthalmoplegia (or nystagmus), and ataxia. However, clinical experience and modern studies have demonstrated that this classic triad is fully present in only a minority of patients, leading to significant underdiagnosis and delayed treatment. Thiamine plays a critical role as a coenzyme in carbohydrate metabolism and cellular energy production, particularly in the brain, which has high metabolic demands. When thiamine levels are depleted, brain regions such as the thalamus, hypothalamus, mammillary bodies, and brainstem experience biochemical failure, leading to cellular swelling, oxidative stress, and eventually, irreversible tissue damage.

While Wernicke encephalopathy has been classically associated with chronic alcohol use disorder, it is increasingly recognized in non-alcoholic populations who suffer from severe malnutrition, gastrointestinal disorders, hyperemesis gravidarum, or prolonged fasting. The condition represents a major medical emergency. Left untreated or inadequately treated, it can progress rapidly to permanent cognitive impairment (Korsakoff syndrome) or result in coma and death. Recognition of its diverse etiologies and prompt administration of intravenous thiamine are crucial in preventing long-term neurological sequelae.

\section*{Epidemiology}
The true prevalence of Wernicke encephalopathy is difficult to estimate because many cases go undiagnosed during life. Autopsy studies suggest a prevalence of 0.8\% to 2.8\% in the general population, with a significantly higher prevalence of up to 12.5\% in individuals with chronic alcohol use disorder. The diagnosis is missed clinically in approximately 80\% of cases, highlighting the need for increased clinical suspicion among emergency physicians, psychiatrists, and general practitioners.

Demographically, Wernicke encephalopathy can occur in individuals of any age, gender, or socioeconomic background, provided there is a persistent thiamine deficiency. However, it is most frequently observed in adults aged 30 to 70 years, with a higher incidence in males, which correlates with the higher prevalence of chronic alcohol consumption in this demographic. In non-alcoholic populations, the female-to-male ratio is more balanced and depends on the underlying etiology, such as hyperemesis gravidarum in pregnant women or eating disorders in young adults. Incidence rates are also elevated in patients undergoing bariatric surgery, those with end-stage renal disease on dialysis, oncological patients undergoing chemotherapy, and populations experiencing severe food insecurity or famine. Geographically, Wernicke encephalopathy occurs worldwide, though the underlying risk factors vary between developed nations (where chronic alcoholism and gastrointestinal surgeries dominate) and developing countries (where nutritional deficiencies and systemic infections are more prevalent).

\section*{Aetiology and Pathophysiology}
Thiamine deficiency leads to Wernicke encephalopathy through a cascade of biochemical events that culminate in cellular dysfunction and cell death in vulnerable brain regions. The etiology and pathophysiology of this condition can be understood through the following mechanisms:

\begin{itemize}
    \item \textbf{Thiamine (Vitamin B1) Deficiency}: Thiamine is a water-soluble vitamin that the human body cannot synthesize endogenously. It has a short half-life, and body stores are typically depleted within 18 to 20 days of zero intake. It is absorbed in the duodenum and proximal jejunum via active transport at low concentrations and passive diffusion at high concentrations.
    \item \textbf{Biochemical Role of Thiamine Pyrophosphate}: Once absorbed, thiamine is phosphorylated into its active form, thiamine pyrophosphate (TPP). TPP serves as an essential coenzyme for several key enzymes in the Krebs cycle and the pentose phosphate pathway, including pyruvate dehydrogenase, alpha-ketoglutarate dehydrogenase, and transketolase. These enzymes are critical for glucose metabolism and the production of adenosine triphosphate (ATP).
    \item \textbf{Energy Failure and Excitotoxicity}: When thiamine levels are deficient, the reduction in TPP-dependent enzyme activity leads to a severe impairment in ATP production. The brain, which relies almost exclusively on glucose for energy, is highly vulnerable. The failure of energy-dependent active transport pumps (such as the Na\textsuperscript{+}/K\textsuperscript{+}-ATPase pump) results in cellular depolarization, intracellular sodium accumulation, and cytotoxic edema.
    \item \textbf{Lactic Acid Accumulation}: The impairment of pyruvate dehydrogenase prevents pyruvate from entering the Krebs cycle, forcing it to convert to lactate. This leads to localized lactic acidosis, which further damages glial cells and neurons.
    \item \textbf{Glutamate-Mediated Neurotoxicity}: Energy depletion prevents the reuptake of the excitatory neurotransmitter glutamate from the synaptic cleft. The resulting excess glutamate overstimulates N-methyl-D-aspartate (NMDA) receptors, leading to an influx of calcium into neurons, activating destructive intracellular enzymes, and causing excitotoxic cell death.
    \item \textbf{Oxidative Stress and Blood-Brain Barrier Disruption}: Thiamine deficiency increases the production of reactive oxygen species (ROS) and reactive nitrogen species, leading to oxidative damage to lipids, proteins, and DNA. It also triggers inflammatory responses, microglial activation, and the disruption of the blood-brain barrier, allowing toxic substances and inflammatory mediators to enter the brain parenchyma.
    \item \textbf{Anatomical Vulnerability}: Specific regions of the brain are highly vulnerable due to their high rate of thiamine turnover and glucose utilization. These include the mammillary bodies, the medial and midline nuclei of the thalamus, the periaqueductal gray matter of the midbrain, the floor of the fourth ventricle, and the cerebellar vermis. Histopathological changes in these areas include congestion of capillaries, petechial hemorrhages, demyelination, axonal damage, and neuronal loss.
    \item \textbf{Risk Factors and Causes}: Common causes of thiamine deficiency include:
    \begin{itemize}
        \item \textbf{Chronic Alcohol Use Disorder}: Alcohol impairs thiamine absorption from the gastrointestinal tract, reduces hepatic storage of the vitamin, impairs its phosphorylation into TPP, and is often accompanied by poor dietary intake.
        \item \textbf{Bariatric and Gastrointestinal Surgeries}: Procedures such as Roux-en-Y gastric bypass, sleeve gastrectomy, or bowel resections can severely limit the mucosal surface area available for thiamine absorption.
        \item \textbf{Severe Gastrointestinal Disease}: Conditions such as Crohn's disease, celiac disease, chronic diarrhea, persistent vomiting, or gastrointestinal malignancies impair nutrient absorption.
        \item \textbf{Hyperemesis Gravidarum}: Prolonged, severe nausea and vomiting during pregnancy can rapidly deplete maternal thiamine stores.
        \item \textbf{Systemic Illnesses and Interventions}: Conditions that increase metabolic demand, such as systemic infections, cancer (especially hematological malignancies), end-stage renal disease on hemodialysis, and prolonged parenteral nutrition without adequate vitamin supplementation, can precipitate Wernicke encephalopathy.
    \end{itemize}
\end{itemize}

\section*{Clinical Features}
The presentation of Wernicke encephalopathy can be highly variable, and the classic triad of signs is present in only 10\% to 16\% of patients. Clinicians must maintain a high index of suspicion, especially in at-risk populations. The primary clinical features include:

\begin{itemize}
    \item \textbf{Encephalopathy (Mental Status Changes)}: This is the most common presenting feature, occurring in approximately 82\% of cases. It typically manifests as global confusion, disorientation to time and place, inattentiveness, and apathy. Patients may display a profound lack of interest in their surroundings, sluggish speech, and lethargy. If left untreated, the encephalopathy can progress to stupor, obtundation, and eventually, coma.
    \item \textbf{Oculomotor Abnormalities}: Ocular signs are present in approximately 29\% of cases and are hallmark indicators of the disease. The most common sign is nystagmus, which can be horizontal, vertical, or rotary. Horizontal nystagmus is typically elicited by lateral gaze. Other oculomotor signs include:
    \begin{itemize}
        \item Abducens nerve (cranial nerve VI) palsy, which is often bilateral and causes weakness or paralysis of the lateral rectus muscles, leading to internal strabismus and diplopia (double vision).
        \item Conjugate gaze palsies, where the patient cannot move both eyes together in a particular direction.
        \item Complete external ophthalmoplegia, where all eye movements are paralyzed.
        \item Pupillary abnormalities, such as sluggish light reflex, anisocoria, or miosis.
    \end{itemize}
    \item \textbf{Ataxia (Gait and Vestibular Dysfunction)}: Ataxia is present in approximately 23\% of patients. It primarily affects the lower extremities, leading to a wide-based, slow, and unsteady gait. In mild cases, the ataxia may only be apparent during tandem walking. In severe cases, the patient may be completely unable to stand or walk without assistance. The ataxia is caused by a combination of cerebellar dysfunction and vestibular hypofunction.
    \item \textbf{Cardiovascular and Autonomic Features}: In some cases, Wernicke encephalopathy can present with signs of cardiovascular compromise, historically referred to as ``wet beriberi''. This includes tachycardia, dyspnea, orthopnea, and peripheral edema due to high-output cardiac failure. Autonomic dysfunction may also manifest as orthostatic hypotension, hypothermia, or unexplained syncope.
    \item \textbf{Peripheral Neuropathy}: Many patients, particularly those with chronic alcohol use disorder, have a coexisting peripheral neuropathy. Symptoms include symmetric numbness, tingling, burning pain, and weakness in the distal limbs (primarily the feet and hands), along with diminished deep tendon reflexes.
\end{itemize}

\section*{Differential Diagnosis}
Due to the non-specific nature of its symptoms, Wernicke encephalopathy can be confused with many other neurological and medical conditions. It is critical to rule these out while simultaneously initiating treatment if suspicion of Wernicke encephalopathy is high. Key differential diagnoses include:

\begin{itemize}
    \item \textbf{Acute Ischemic or Hemorrhagic Stroke}: Particularly strokes involving the brainstem, thalamus, or cerebellum, which can present with sudden onset of ataxia, oculomotor deficits, and altered mental status.
    \item \textbf{Sepsis and Severe Infection}: Systemic infections, meningitis, or encephalitis can cause acute encephalopathy, delirium, and confusion, resembling the mental status changes of Wernicke encephalopathy.
    \item \textbf{Hypoglycemia}: Low blood glucose levels can mimic the acute mental status changes, confusion, and focal neurological signs of Wernicke encephalopathy and must be ruled out immediately with a finger-prick test.
    \item \textbf{Hepatic Encephalopathy}: Patients with chronic liver disease (frequently comorbid with alcohol use disorder) can present with confusion, asterixis (flapping tremor), and altered consciousness due to elevated blood ammonia levels.
    \item \textbf{Alcohol Withdrawal and Delirium Tremens}: Alcohol withdrawal can cause agitation, confusion, hallucinations, autonomic instability, and tremors. While it can coexist with Wernicke encephalopathy, it requires different management (e.g., benzodiazepines).
    \item \textbf{Electrolyte Imbalances}: Severe hyponatremia, hypernatremia, hypokalemia, or hypercalcemia can cause profound encephalopathy, weakness, and confusion.
    \item \textbf{Sedative-Hypnotic Intoxication or Withdrawal}: Overdose or withdrawal from benzodiazepines, barbiturates, or other sedatives can cause ataxia, slurred speech, and depressed mental status.
    \item \textbf{Dementia (e.g., Alzheimer's disease or Vascular Dementia)}: Chronic cognitive decline can be confused with the encephalopathy of Wernicke, though dementia is typically slower in onset and progressive.
    \item \textbf{Non-convulsive Status Epilepticus}: Ongoing, subclinical seizure activity can present as unexplained confusion, responsiveness fluctuations, or stupor.
    \item \textbf{Hypoxia and Carbon Monoxide Poisoning}: Deprivation of oxygen to the brain can cause acute cognitive deficits, confusion, and motor instability.
\end{itemize}

\section*{When to Seek Medical Care}
Wernicke encephalopathy is a medical emergency that requires immediate intervention. If an individual displays any symptoms of acute confusion, difficulty walking, or abnormal eye movements, emergency medical services must be contacted immediately.

In life-threatening situations, dial the local emergency number:
\begin{itemize}
    \item \textbf{local emergency services} in many countries.
    \item \textbf{911} in the United States and Canada.
    \item \textbf{999} in the United Kingdom.
    \item \textbf{112} in the European Union and many other countries.
\end{itemize}

Do not wait for all symptoms of the classic triad to appear before seeking help. Delayed treatment can lead to permanent brain damage, coma, or death. Inform emergency responders if the patient has a history of heavy alcohol consumption, recent bariatric surgery, severe weight loss, chronic vomiting, or poor nutrition, as this information is critical for guiding immediate treatment in the emergency department.

\section*{Diagnosis and Investigations}
The diagnosis of Wernicke encephalopathy is primarily clinical. There is no single diagnostic test that can definitively confirm the condition. Instead, diagnostic criteria, physical examinations, and investigations are used to support the clinical diagnosis and exclude other causes.

\begin{itemize}
    \item \textbf{Clinical Diagnostic Criteria (Caine Criteria)}: The Caine criteria are widely used to identify patients at risk of Wernicke encephalopathy. A diagnosis is made if a patient satisfies at least two of the following four criteria:
    \begin{enumerate}
        \item Dietary deficiency (e.g., severe malnutrition, alcohol abuse, chronic vomiting).
        \item Ocular signs (nystagmus, ophthalmoplegia, bilateral abducens palsy).
        \item Cerebellar dysfunction (ataxia of gait, unsteady posture).
        \item Altered mental status or mild memory impairment.
    \end{enumerate}
    \item \textbf{Laboratory Investigations}:
    \begin{itemize}
        \item \textbf{Blood Thiamine Levels}: Measurement of whole-blood thiamine diphosphate levels using high-performance liquid chromatography (HPLC) can help assess thiamine status, but results are rarely available in real-time to guide acute management. Treatment must not be delayed while waiting for these results.
        \item \textbf{Erythrocyte Transketolase Activity (ETKA)}: This assay measures the activity of the transketolase enzyme in red blood cells before and after the addition of thiamine pyrophosphate. An increase in enzyme activity of more than 15\% to 25\% suggests thiamine deficiency. However, this test is technically difficult, rarely available, and lacks specificity.
        \item \textbf{General Laboratory Panel}: Complete blood count (to check for anemia), blood glucose (to rule out hypoglycemia), liver function tests (to evaluate for hepatic encephalopathy), renal function, and serum electrolytes (sodium, potassium, magnesium, calcium) are essential. Magnesium levels are particularly important, as magnesium is a cofactor for thiamine-dependent enzymes, and hypomagnesemia can render thiamine treatment ineffective.
    \end{itemize}
    \item \textbf{Neuroimaging}:
    \begin{itemize}
        \item \textbf{Magnetic Resonance Imaging (MRI)}: MRI is the imaging modality of choice. It has a sensitivity of approximately 53\% and a specificity of 93\%. Characteristic findings on T2-weighted and fluid-attenuated inversion recovery (FLAIR) sequences include symmetrical hyperintensities in the medial thalami, periaqueductal gray matter, tectal plate, and mammillary bodies. Mammillary body atrophy and contrast enhancement are highly specific features, especially in chronic or alcoholic cases.
        \item \textbf{Computed Tomography (CT) Scan}: A head CT is generally insensitive for Wernicke encephalopathy and is primarily used to exclude other acute intracranial pathologies, such as hemorrhage, mass lesions, or large strokes.
    \end{itemize}
    \item \textbf{Lumbar Puncture (Cerebrospinal Fluid Analysis)}: A lumbar puncture is not routinely indicated for diagnosing Wernicke encephalopathy but may be performed to rule out central nervous system infections (meningitis, encephalitis) or inflammatory conditions if the presentation is atypical.
\end{itemize}

\section*{Management}
The management of Wernicke encephalopathy requires urgent administration of parenteral thiamine. Oral thiamine is highly unreliable in the acute phase due to impaired gastrointestinal absorption, particularly in patients with alcohol use disorder or gastrointestinal disease.

\begin{itemize}
    \item \textbf{Intravenous Thiamine Administration}:
    \begin{itemize}
        \item \textbf{Dosage and Frequency}: Use the local hospital or specialist protocol; commonly, high-dose intravenous thiamine is given three times daily for several days, followed by lower-dose parenteral treatment until improvement. Exact dose, duration, route, and dilution vary by guideline and patient factors, so treatment must not be self-prescribed.
        \item \textbf{Safety and Anaphylaxis}: Intravenous thiamine is generally very safe. Although rare cases of anaphylaxis have been reported, the risk is extremely low, and the benefits of preventing irreversible brain damage far outweigh the risks.
    \end{itemize}
    \item \textbf{Glucose and Thiamine}: Give parenteral thiamine promptly when deficiency is suspected, preferably before or alongside carbohydrate when feasible. Do not delay life-saving glucose for severe hypoglycaemia; give glucose immediately and thiamine concurrently or as soon as possible.
    \item \textbf{Magnesium Supplementation}: Magnesium is a crucial cofactor for many thiamine-dependent enzymes. In patients with Wernicke encephalopathy, particularly those with a history of alcohol use disorder, hypomagnesemia is common. If magnesium levels are low, thiamine may be therapeutically ineffective. Serum magnesium levels should be corrected using intravenous magnesium sulfate (typically 1 to 2~grams) alongside thiamine replacement.
    \item \textbf{Correction of Electrolytes and Hydration}: Intravenous fluids should be administered to maintain hydration, using normal saline or balanced crystalloid solutions. Other electrolyte imbalances (such as hypokalemia or hypophosphatemia) must be carefully monitored and corrected.
    \item \textbf{Nutritional Support and Rehabilitation}: Once the acute phase is stabilized and gastrointestinal function is restored, the patient should transition to oral thiamine supplementation (typically 100~mg to 300~mg daily) and a balanced diet. Nutritional assessment and rehabilitation therapy (physical and occupational therapy) are crucial for recovering motor skills, balance, and activities of daily living.
\end{itemize}

\section*{Complications}
Untreated or inadequately treated Wernicke encephalopathy carries a very high morbidity and mortality rate. The primary complications include:

\begin{itemize}
    \item \textbf{Korsakoff Syndrome (Korsakoff Amnestic State)}: This is the most common chronic complication of Wernicke encephalopathy, occurring in up to 80\% of untreated survivors. It is characterized by severe anterograde amnesia (inability to form new memories) and retrograde amnesia (loss of past memories). Patients often exhibit ``confabulation'', where they unconsciously fabricate stories to fill in memory gaps. They may also display apathy, loss of initiative, and cognitive deficits. When Wernicke encephalopathy and Korsakoff syndrome occur sequentially or together, the condition is referred to as Wernicke-Korsakoff syndrome (WKS).
    \item \textbf{Permanent Neurological Deficits}: Even after treatment, some patients may suffer from persistent neurological impairments, such as permanent ataxia, requiring assistive devices for mobility, chronic nystagmus, or permanent oculomotor weakness.
    \item \textbf{Coma and Death}: Severe, untreated thiamine deficiency leads to progressive energy failure in critical brainstem centers regulating respiration and cardiovascular function. This can lead to obtundation, coma, and ultimately, death. The mortality rate of untreated Wernicke encephalopathy is estimated to be approximately 10\% to 20\%.
    \item \textbf{Cardiovascular Failure (Wet Beriberi)}: Severe thiamine deficiency can lead to high-output cardiac failure, pulmonary edema, and cardiovascular collapse.
    \item \textbf{Treatment-Related Complications}: Rapid correction of electrolyte abnormalities, particularly hyponatremia, carries a risk of osmotic demyelination syndrome if not managed carefully. Additionally, volume overload can occur in patients with compromised cardiac or renal function receiving intensive intravenous fluid therapy.
\end{itemize}

\section*{Prognosis and Outcomes}
The prognosis of Wernicke encephalopathy depends heavily on the speed of diagnosis and the promptness of parenteral thiamine administration. If treated early, many of the neurological signs are reversible:

Ocular abnormalities, such as abducens nerve palsy and conjugate gaze deficits, typically show the most rapid improvement, often resolving within hours to days after the first thiamine doses. Nystagmus may persist for weeks or, in some cases, become permanent. The acute encephalopathy and confusion resolve gradually over days to weeks, although complete cognitive recovery may take months. Ataxia is the slowest symptom to improve. Approximately 35\% to 40\% of patients recover completely, but the remainder may have a residual wide-based gait or require walking aids.

If treatment is delayed or inadequate, the prognosis is much poorer. About 80\% of patients who survive the acute phase of Wernicke encephalopathy will develop Korsakoff syndrome. Of these, only about 20\% achieve complete recovery over time, 50\% show partial improvement, and the remaining 30\% require long-term institutional care due to severe, irreversible memory loss. The presence of comorbid conditions, such as advanced liver disease, severe malnutrition, or systemic infections, significantly worsens the prognosis.

\section*{Prevention and Self-Care}
Preventing Wernicke encephalopathy involves addressing the underlying risk factors for thiamine deficiency, particularly in high-risk populations. Key preventive strategies and self-care measures include:

\begin{itemize}
    \item \textbf{Nutritional Fortification and Diet}: Consuming a diet rich in thiamine is the most effective preventative measure. Foods high in thiamine include whole grains, fortified cereals, legumes (beans, lentils), pork, yeast extract, nuts, and seeds. Many countries mandate the fortification of wheat flour and cereals with thiamine, which has significantly reduced the incidence of nutritional deficiency.
    \item \textbf{Prophylaxis in High-Risk Patients}: Patients admitted to hospitals with a history of chronic alcohol use disorder, severe malnutrition, or persistent vomiting should receive prophylactic parenteral thiamine (e.g., 100~mg to 250~mg daily) even in the absence of neurological symptoms, particularly before administering any glucose-containing fluids.
    \item \textbf{Care After Bariatric Surgery}: Individuals who have undergone bariatric surgery must adhere to lifelong vitamin supplementation guidelines, including daily multivitamin formulas containing thiamine, and should receive periodic blood tests to monitor nutritional levels.
    \item \textbf{Alcohol Cessation and Rehabilitation}: For individuals with alcohol use disorder, achieving and maintaining sobriety is crucial. Alcohol rehabilitation programs, counseling, support groups (such as Alcoholics Anonymous), and medical therapies to reduce alcohol cravings can prevent recurrent episodes of thiamine deficiency.
    \item \textbf{Management of Hyperemesis Gravidarum}: Pregnant women suffering from severe, persistent vomiting should be monitored closely and treated with antiemetics, intravenous hydration, and early thiamine supplementation to prevent rapid depletion of body stores.
    \item \textbf{Self-Monitoring and Early Reporting}: Patients at risk of nutritional deficiencies, along with their families and caregivers, should be educated on the early warning signs of thiamine deficiency, such as mild confusion, unsteady gait, or double vision, and instructed to seek immediate medical attention if these symptoms arise.
\end{itemize}

\section*{Sources and Further Reading}
\begin{itemize}
    \item \textbf{World Health Organization (WHO)}: Nutritional Deficiencies and Global Health Guidelines. Website: \texttt{https://www.who.int/}
    \item \textbf{National Institute for Health and Care Excellence (NICE)}: Alcohol-use disorders: diagnosis and management of physical complications. Website: \texttt{https://www.nice.org.uk/}
    \item \textbf{National Institutes of Health (NIH) - Office of Dietary Supplements}: Thiamine Fact Sheet for Health Professionals. Website: \texttt{https://ods.od.nih.gov/}
    \item \textbf{MedlinePlus (National Library of Medicine)}: Wernicke-Korsakoff Syndrome Medical Information. Website: \texttt{https://medlineplus.gov/}
    \item \textbf{National Institute on Alcohol Abuse and Alcoholism (NIAAA)}: Neurological Effects of Alcohol and Thiamine Deficiency. Website: \texttt{https://www.niaaa.nih.gov/}
    \item \textbf{UpToDate}: Clinical Manifestations, Diagnosis, and Treatment of Wernicke Encephalopathy. Website: \texttt{https://www.uptodate.com/}
\end{itemize}
